% ADDED figure (2026-11-13): Theorem §C6c.9.5, the Feynman rules of the broken phase of the linear sigma model (PHY 513, Problem Set 10, Problem 1(b)): two propagators, two cubic and three quartic vertices; pion lines solid blue, sigma lines dashed red. Source kept in qft-C/tools/hw10/figs/.
\tikzset{
  fpi/.style={line width=0.9pt, figB},
  fsg/.style={line width=0.9pt, figR, densely dashed},
  cv/.style={circle, fill=figO, draw=figO, inner sep=0pt, minimum size=5pt},
  fdlab/.style={font=\scriptsize, inner sep=2.2pt},
  fdS/.style={font=\small},
}
\begin{tikzpicture}[font=\small]
  % propagators
  \begin{scope}[xshift=0cm, yshift=0cm]
    \draw[fpi] (-0.9,0.45)--(0.9,0.45); \node[fdlab,left] at (-0.9,0.45){$i$}; \node[fdlab,right] at (0.9,0.45){$j$}; \node[fdlab,above] at (0,0.45){$p$};
    \node[fdS] at (0,-0.25) {$\pi$:\ \ $\dfrac{i\,\delta^{ij}}{p^2 + i\epsilon}$};
    \draw[fsg] (-0.9,-1.45)--(0.9,-1.45); \node[fdlab,above] at (0,-1.45){$p$};
    \node[fdS] at (0,-2.15) {$\sigma$:\ \ $\dfrac{i}{p^2 - 2\mu^2 + i\epsilon}$};
  \end{scope}
  % pi pi sigma
  \begin{scope}[xshift=3.6cm, yshift=0.2cm]
    \draw[fpi] (-0.8,0.7)--(0,0); \draw[fpi] (-0.8,-0.7)--(0,0); \draw[fsg] (0,0)--(0.95,0);
    \node[cv] at (0,0){};
    \node[fdlab,left] at (-0.8,0.7){$i$}; \node[fdlab,left] at (-0.8,-0.7){$j$};
    \node[fdS] at (0.05,-1.25) {$-2i\lambda v\,\delta^{ij}$};
  \end{scope}
  % sigma^3
  \begin{scope}[xshift=3.6cm, yshift=-2.75cm]
    \draw[fsg] (-0.8,0.7)--(0,0); \draw[fsg] (-0.8,-0.7)--(0,0); \draw[fsg] (0,0)--(0.95,0);
    \node[cv] at (0,0){};
    \node[fdS] at (0.05,-1.25) {$-6i\lambda v$};
  \end{scope}
  % pi pi sigma sigma
  \begin{scope}[xshift=7.4cm, yshift=0.2cm]
    \draw[fpi] (-0.8,0.7)--(0,0); \draw[fpi] (-0.8,-0.7)--(0,0); \draw[fsg] (0,0)--(0.8,0.7); \draw[fsg] (0,0)--(0.8,-0.7);
    \node[cv] at (0,0){};
    \node[fdlab,left] at (-0.8,0.7){$i$}; \node[fdlab,left] at (-0.8,-0.7){$j$};
    \node[fdS] at (0,-1.25) {$-2i\lambda\,\delta^{ij}$};
  \end{scope}
  % sigma^4
  \begin{scope}[xshift=10.6cm, yshift=0.2cm]
    \draw[fsg] (-0.8,0.7)--(0,0); \draw[fsg] (-0.8,-0.7)--(0,0); \draw[fsg] (0,0)--(0.8,0.7); \draw[fsg] (0,0)--(0.8,-0.7);
    \node[cv] at (0,0){};
    \node[fdS] at (0,-1.25) {$-6i\lambda$};
  \end{scope}
  % pi^4
  \begin{scope}[xshift=9.0cm, yshift=-2.75cm]
    \draw[fpi] (-0.8,0.7)--(0,0); \draw[fpi] (-0.8,-0.7)--(0,0); \draw[fpi] (0,0)--(0.8,0.7); \draw[fpi] (0,0)--(0.8,-0.7);
    \node[cv] at (0,0){};
    \node[fdlab,left] at (-0.8,0.7){$i$}; \node[fdlab,left] at (-0.8,-0.7){$j$};
    \node[fdlab,right] at (0.8,0.7){$k$}; \node[fdlab,right] at (0.8,-0.7){$l$};
    \node[fdS] at (0,-1.3) {$-2i\lambda\,\bigl(\delta^{ij}\delta^{kl} + \delta^{ik}\delta^{jl} + \delta^{il}\delta^{jk}\bigr)$};
  \end{scope}
  \node[fdS] at (0,1.55) {propagators};
  \node[fdS] at (3.6,1.55) {cubic vertices};
  \node[fdS] at (9.0,1.55) {quartic vertices};
  \draw[gray!50, line width=0.3pt] (1.9,1.75)--(1.9,-4.4); \draw[gray!50, line width=0.3pt] (5.5,1.75)--(5.5,-4.4);
\end{tikzpicture}
